% Generated draft: requires mathematical and source review.
\documentclass[11pt]{article}
\usepackage[margin=25mm]{geometry}
\usepackage{fontspec}
\setmainfont{Latin Modern Roman}
\usepackage{amsmath,amssymb,mathtools,unicode-math}
\setmathfont{Latin Modern Math}
\usepackage{enumitem,needspace,xurl,hyperref}
\setlength{\parindent}{0pt}
\setlength{\parskip}{5pt}
\setlength{\emergencystretch}{3em}
\newcommand{\N}{\mathbb{N}}
\newcommand{\Z}{\mathbb{Z}}
\newcommand{\Q}{\mathbb{Q}}
\newcommand{\R}{\mathbb{R}}
\newcommand{\C}{\mathbb{C}}
\newcommand{\F}{\mathbb{F}}
\newcommand{\A}{\mathbb{A}}
\newcommand{\PP}{\mathbb{P}}
\newcommand{\E}{\mathbb{E}}
\newcommand{\eps}{\varepsilon}
\newcommand{\dd}{\,\mathrm{d}}
\begin{document}
\section*{UnsolvedMath ID: 1101213}
\textbf{Source problem code:} catalogue entry
\subsection{Definitions and mathematical statement}
Let $S_g$ be a closed connected orientable surface of genus $g\geq2$, and let $\operatorname{Mod}(S_g)$ be the group of isotopy classes of orientation-preserving homeomorphisms of $S_g$. A mapping class is pseudo-Anosov if it has a representative preserving a pair of transverse measured singular foliations and scaling their transverse measures by $\lambda$ and $\lambda^{-1}$ for some $\lambda>1$. A subgroup $H\leq\operatorname{Mod}(S_g)$ is purely pseudo-Anosov if every nonidentity element is pseudo-Anosov. It is non-free if it is not isomorphic to a free group.

The question is whether a non-free purely pseudo-Anosov subgroup exists. The original isolated wording does not specify the surface class; the supplied resolving assessment gives the following sufficient closed-surface setting:
\[
 \exists g\geq2\ \exists H\leq\operatorname{Mod}(S_g),\qquad
 H\text{ is non-free and purely pseudo-Anosov}.
\]
The identity is excluded from the pseudo-Anosov requirement, as is necessary in the usual subgroup terminology.

The supplied assessment reports the stronger existence statement that, for every $g\geq4$, there are purely pseudo-Anosov subgroups of $\operatorname{Mod}(S_g)$ isomorphic to fundamental groups of closed surfaces of genus at least two, and infinitely many commensurability classes of such subgroups. The sense of commensurability is not specified in the record. In an ambient group, subgroups may be called commensurable if their intersection has finite index in both; another convention allows conjugating one subgroup first. Abstract commensurability of the surface groups themselves is a different notion and cannot distinguish closed hyperbolic surface groups in the way the summary claims. The stronger multiplicity assertion is retained with this unresolved convention rather than silently choosing one.
\subsection{Short English statement}
Can a subgroup of a mapping class group be non-free while every nonidentity element is pseudo-Anosov? The supplied assessment answers yes using closed surface subgroups, attributes the result to Kent--Leininger (2024), and labels the problem solved in the literature. It further reports infinitely many commensurability classes in every closed genus at least four and a construction of the first closed aspherical atoroidal surface bundles over surfaces. Those are attributed literature claims. The assessment's statement that it checked the arXiv record, author list and journal status does not constitute a source visit or independent proof by this draft.
\subsection{Sources}
\begin{itemize}
\item UnsolvedMath permanent ID 1101213, source code catalogue entry, ``Questions in Geometric Group Theory --- Q 12.13.'' Registry curation and metadata: Ulam AI and UnsolvedMath Contributors, CC BY 4.0; updated 2026-08-06.
\item Bestvina, \emph{Questions in Geometric Group Theory} (source list dated 2004), Question 12.13, PDF page 22: \url{https://www.math.utah.edu/~bestvina/eprints/questions-updated.pdf}. Supplied manual transcription accessed 2026-07-29. Original metadata presents an open item without independent current-status verification; status/release and rights review are NEEDS\_REVIEW.
\item Kent--Leininger, \emph{Atoroidal surface bundles}, arXiv:2405.12067, \url{https://arxiv.org/abs/2405.12067}, described by the supplied assessment as to appear in \emph{Annals of Mathematics}. This publication-status description is retained as supplied, without asserting a later update.
\end{itemize}
\end{document}
